% 内积空间：由内积定义的长度、正交、Cauchy--Schwarz不等式与夹角。

\section{由内积定义的夹角}
\label{sec:ls-inner-product-angle}

\paragraph{长度与正交}
与$\mathbb R^d$中一样，内积给出长度$\lVert\bm v\rVert=\sqrt{\langle\bm v,\bm v\rangle}$，它是否满足范数的全部要求，留到下一节检验；$\langle\bm x,\bm y\rangle=0$时称$\bm x,\bm y$正交。正交的两个向量满足勾股定理。展开
\[
 \lVert\bm x+\bm y\rVert^2=\lVert\bm x\rVert^2+2\langle\bm x,\bm y\rangle+\lVert\bm y\rVert^2,
\]
交叉项为零，得$\lVert\bm x+\bm y\rVert^2=\lVert\bm x\rVert^2+\lVert\bm y\rVert^2$。

\paragraph{Cauchy--Schwarz不等式}
平面上两个非零向量的夹角$\theta$满足
\[
 \langle\bm x,\bm y\rangle=\lVert\bm x\rVert\,\lVert\bm y\rVert\cos\theta.
\]
要在一般内积空间中反过来用这个关系定义夹角，须先保证内积除以两个长度之积后总落在$[-1,1]$中。下面的不等式给出这一保证。
\begin{theorem}[Cauchy--Schwarz不等式]
\label{thm:linear-algebra-cauchy-schwarz}
对任一内积及其长度，任意向量$\bm x,\bm y$满足
\begin{equation}
 |\langle\bm x,\bm y\rangle|\leq\lVert\bm x\rVert\,\lVert\bm y\rVert,
 \label{eq:prelim-cauchy-schwarz}
\end{equation}
等号成立当且仅当其中一个向量是另一个向量的实数倍。
\end{theorem}
\begin{proof}
$\bm y=\bm0$时两边都为零。设$\bm y\neq\bm0$。对任意实数$s$，向量$\bm x-s\bm y$的长度平方非负；按内积的对称性与双线性展开，
\[
 0\leq\lVert\bm x-s\bm y\rVert^2
 =\lVert\bm x\rVert^2-2s\langle\bm x,\bm y\rangle+s^2\lVert\bm y\rVert^2.
\]
右侧是$s$的二次函数，二次项系数为正，在顶点$s=\langle\bm x,\bm y\rangle/\lVert\bm y\rVert^2$处最小，最小值为$\lVert\bm x\rVert^2-\langle\bm x,\bm y\rangle^2/\lVert\bm y\rVert^2$。它非负，两边乘以$\lVert\bm y\rVert^2$再开平方，即得不等式。等号成立当且仅当对这个$s$有$\bm x-s\bm y=\bm0$，即$\bm x$是$\bm y$的实数倍；连同$\bm y=\bm0$的情形，即得等号条件。
\end{proof}
证明只用了内积的三条公理，所以不等式对每一种内积都成立。在$\mathbb R^d$的标准内积下，它就是$p=q=2$的Hölder不等式（定理~\ref{thm:linear-algebra-holder}）；对Frobenius内积，它给出$|\langle\bm A,\bm B\rangle_{\mathrm F}|\leq\lVert\bm A\rVert_{\mathrm F}\lVert\bm B\rVert_{\mathrm F}$。

\paragraph{夹角}
由Cauchy--Schwarz不等式，非零向量的夹角$\theta\in[0,\pi]$可以由
\[
 \cos\theta=\frac{\langle\bm x,\bm y\rangle}{\lVert\bm x\rVert\,\lVert\bm y\rVert}
 =\left\langle\frac{\bm x}{\lVert\bm x\rVert},\frac{\bm y}{\lVert\bm y\rVert}\right\rangle
\]
定义。第二种写法说明，夹角只比较方向：先把两个向量各自归一化为单位向量，再取内积，所以$\bm x$与$c\bm x$（$c>0$）给出同一个方向。正交恰好对应$\theta=\pi/2$，这就是把正交称为垂直的理由。由等号条件，$\cos\theta=\pm1$恰好对应两个向量共线。分子和分母必须来自同一个内积；一边用$\ell_2$、一边用$\ell_1$计量长度时，这个比值没有几何意义。

回到第~\ref{chap:linear-operations}章贯穿例中两个特征的取值向量，计算它们的夹角。
\begin{example}[贯穿例两列的夹角]
\label{ex:ls-nearly-collinear-angle}
贯穿例的两列是$\bm c_1=(1,1,1)^\top$与$\bm c_2=(1,1+\varepsilon,1-\varepsilon)^\top$，取欧氏内积。夹角公式需要一个内积和两个长度：$\langle\bm c_1,\bm c_2\rangle=1+(1+\varepsilon)+(1-\varepsilon)=3$，其中$\pm\varepsilon$恰好抵消；$\lVert\bm c_1\rVert_2^2=3$；$\lVert\bm c_2\rVert_2^2=1+(1+\varepsilon)^2+(1-\varepsilon)^2=3+2\varepsilon^2$，一次项同样抵消，只剩下二次项。所以
\begin{equation}
 \cos\theta_\varepsilon=\frac{3}{\sqrt3\,\sqrt{3+2\varepsilon^2}}
 =\frac{1}{\sqrt{1+2\varepsilon^2/3}}\longrightarrow1\qquad(\varepsilon\to0).
 \label{eq:linear-algebra-nearly-collinear-angle}
\end{equation}
代入具体数值：$\varepsilon=0.1$时$\cos\theta_\varepsilon\approx0.9967$，夹角约$4.7^\circ$；$\varepsilon=0.01$时$\cos\theta_\varepsilon\approx0.99997$，夹角不到$0.5^\circ$。

可见$\varepsilon\to0$时两列的夹角趋于零。$\varepsilon$非零时两个向量并不重合，但$\varepsilon$很小时它们几乎同向：$\varepsilon$只以平方出现在余弦中，余弦与$1$只差$\varepsilon^2$量级。
\end{example}
两个向量不重合却几乎同向，称为\term{近共线}（nearly collinear）。这里的“近”指两列作为$\mathbb R^3$中向量的夹角很小，不是说各个样本点彼此接近。夹角也依赖所用的内积：换一种内积，同样两个向量的夹角会改变。
